%% kapa: exam=2 round=26 year=2015 period=2 subject=appraisal_theory subject_ko=감정평가이론 seq=1 number=1 points=40
%% source: Q-Net 감정평가사 기출문제 (cst00309) artlSeq=5149601 file=extracted/1008/5149601/2015년 제26회 감정평가사 2차시험 기출문제/2교시(감정평가이론).pdf
%% conversion: kordoc 4.24.51 -> tools/gen2.py (auto); review=proofread-gpt-sol; flags=none
\begin{longquestion}{【문제 1】}
A법인은 토지 200 m\textsuperscript{2} 및 위 지상에 건축된 연면적 100 m\textsuperscript{2} 1층 업무용 건물(집합건물이 아님)을 소유하고 있다. 건물은 101호 및 102호로 구획되어 있으며, 101호는 A법인이 사무실로 사용하고 있고 102호는 B에게 임대하고 있다. 다음 물음에 답하시오. (40점) \par
1) A법인이 소유한 위 부동산(토지 및 건물)을 감정평가 할 경우 감정평가 규칙에 따른 원칙적인 감정평가방법 및 근거, 해당 방법의 적정성을 논하시오. (15점) \par
2) 임차인 C가 101호를 전세로 임차하기로 하였다. C는 전세금액 및 전세권 설정에 참고하기 위하여 101호 건물 50 m\textsuperscript{2}만을 감정평가 의뢰하였다. 본건 평가의 타당성에 관해 설명하시오. (10점) \par
3) A법인은 토지에 저당권을 설정한 이후 건물을 신축하였으나 건물에 대해서는 저당권을 설정하지 않았다. A법인이 이자지급을 연체하자 저당권자가 본건 토지의 임의경매를 신청하였다. 이 경우 토지의 감정평가 방법에 관해 설명하시오. (5점) \par
4) 해당 토지의 용적률은 50 \%이나 주변토지의 용적률은 100 \%이다. A법인이 용적률 100 \%를 조건으로 하는 감정평가를 의뢰하였다. 조건부평가에 관해 설명하고 본건의 평가 가능 여부를 검토하시오. (10점) \par
\end{longquestion}
